\section{REFERENCIAL TEÓRICO}

\subsection{Avaliações online: compressão das notas e vieses de relato}

Sistemas de avaliação produzem notas concentradas nos valores mais altos. \citet{hu2009} descrevem a distribuição em forma de J das avaliações de produtos, com muitas notas máximas, algumas mínimas e poucas intermediárias e atribuem o padrão a vieses de quem decide avaliar. \citet{dellarocas2008} mostram, com dados do eBay, que a propensão a relatar depende do resultado da transação, o que distorce a distribuição das avaliações públicas em relação à experiência efetivamente vivida. \citet{filippas2022} tratam a inflação de reputação, isto é, a elevação das notas ao longo do tempo, como fator que reduz o conteúdo informativo desses sistemas no longo prazo. O efeito das avaliações sobre as vendas de livros é documentado há tempo \citep{chevalier2006} e as respostas das empresas às avaliações também alteram o que os clientes escrevem \citep{proserpio2017}.

Na psicometria, a tendência de um julgamento global contaminar os julgamentos específicos recebe o nome de efeito halo \citep{thorndike1920}. A noção importa para este estudo porque as subnotas oficiais do Airbnb (limpeza, localização, check-in e outras) são muito correlacionadas entre si. A correlação entre o sentimento de um aspecto e a sua subnota precisa, por isso, ser lida à luz desse efeito.

\subsection{As avaliações do Airbnb e o seu conteúdo}

No Airbnb, \citet{zervas2021} reportam que a grande maioria dos anúncios tem média de 4,5 ou 5 estrelas, valores mais altos que os de hotéis e pousadas listados no TripAdvisor. O sistema é bilateral, pois anfitrião e hóspede avaliam um ao outro. \citet{fradkin2021} realizaram um experimento na plataforma em que o feedback ficava oculto até que as duas partes avaliassem e observaram menos retaliação e reciprocidade nas avaliações e, em consequência, notas mais baixas. \citet{bridges2018} analisam a linguagem das avaliações e concluem que elas recorrem a um conjunto restrito de recursos linguísticos, o que estabelece uma norma de comentário muito positivo, embora experiências menos favoráveis possam ser sinalizadas por indícios mais sutis. Um trabalho com dados do Brasil e dos Estados Unidos compara o sentimento das avaliações do Airbnb, do Couchsurfing e do Booking.com e encontra avaliações consideravelmente mais positivas nas plataformas de economia compartilhada \citep{santos2020}.

Diversos estudos aplicam mineração de texto às avaliações do Airbnb. \citet{cheng2019} concluem que localização, comodidades e anfitrião são os atributos que mais moldam a experiência do hóspede e que o preço não figura entre os determinantes principais. \citet{luo2019} adaptam a análise latente de aspectos e notas para identificar cinco aspectos (comunicação, experiência, localização, produto ou serviço e valor) e mostrar sua contribuição à nota numérica. \citet{sutherland2020} extraem 43 tópicos de mais de um milhão de avaliações de Nova York com LDA, agrupados em avaliação geral do hóspede, atributos de localização, características dos imóveis, gestão do anúncio e qualidade do serviço do anfitrião. Este trabalho se soma a essa linha ao decompor a nota por aspecto em um mercado da América Latina, em três idiomas e com validação explícita da medida.

\subsection{Determinantes de preço e de reputação}

Outra vertente da literatura procura explicar o preço e a demanda dos anúncios. Estudos de preço com regressão hedônica, quantílica e espacial mostram que localização, tamanho, comodidades e reputação online se associam ao preço \citep{wang2017,chicaolmo2020,gyodi2021}. Em Nova York, a localização, sobretudo a acessibilidade ao transporte, tem grande efeito sobre o preço e a receita. A profissionalização do mercado concentra a receita em uma fração menor de anfitriões \citep{deboosere2019}. Esses trabalhos tomam preço, ocupação ou receita como variáveis de interesse. Neste estudo o objeto é a nota e o texto que a acompanha. As medidas de volume entram como desfecho auxiliar, com a ressalva, estabelecida na Seção 3, de que os campos de ocupação e receita estimadas da base são funções do número de avaliações.

\subsection{Análise de sentimento por aspecto e modelos de linguagem}

A análise de sentimento por aspecto consiste em identificar os aspectos de uma entidade mencionados no texto e o sentimento expresso sobre cada um \citep{liu2012}. As tarefas compartilhadas do SemEval consolidaram a formulação e os conjuntos de dados da área \citep{pontiki2014,pontiki2016} e levantamentos recentes organizam os métodos e os desafios \citep{schouten2016,zhang2023absa}. Os modelos atuais de linguagem derivam do BERT \citep{devlin2019}. Para corpora multilíngues como o desta pesquisa, empregam-se modelos multilíngues como o XLM-T, treinado em mensagens de uma rede social em vários idiomas \citep{barbieri2022}. Para o português, existem modelos específicos, como o BERTimbau \citep{souza2020} e trabalhos de mineração de opinião sobre avaliações em português \citep{freitas2013}. Não identificamos, na busca realizada, estudo anterior de análise de sentimento por aspecto sobre avaliações de hospedagem do Rio de Janeiro.

\subsection{Modelagem de tópicos}

O LDA \citep{blei2003} é o método clássico de modelagem de tópicos probabilística. O BERTopic \citep{grootendorst2022} agrupa representações vetoriais de sentenças \citep{reimers2019,reimers2020}, após redução de dimensão por UMAP \citep{mcinnes2018umap} e agrupamento por densidade \citep{mcinnes2017} e descreve cada grupo por termos de maior peso. A coerência $C_v$ \citep{roder2015} é a medida usual para comparar a qualidade interpretável dos tópicos e neste estudo é empregada apenas de forma descritiva.

\subsection{Validação, concordância e importância relativa}

A concordância entre avaliadores é medida pelo coeficiente kappa \citep{cohen1960}, em sua versão ponderada quando as categorias são ordenadas \citep{cohen1968} e sua interpretação segue faixas convencionais \citep{landis1977}, com as ressalvas discutidas em \citet{mchugh2012}. Quando a amostra rotulada é estratificada e se deseja estimar a medida na população, cada observação recebe o peso do seu estrato e o tamanho efetivo da amostra ponderada, descrito por \citet{kish1992}, indica a precisão efetiva. A incerteza é estimada por \textit{bootstrap} \citep{efron1979}. A confiabilidade de uma medida obtida de duas metades dos dados é ajustada pela fórmula de Spearman-Brown \citep{brown1910,spearman1910}. A importância relativa de preditores correlacionados é decomposta pelo método LMG, equivalente à média de Shapley do $R^2$ \citep{lindeman1980,gromping2006,gromping2007}. A correção de Benjamini e Hochberg \citep{benjamini1995} controla a taxa de falsas descobertas em comparações múltiplas e o encolhimento bayesiano empírico \citep{efron1975} fundamenta a opção de ponderar a média de anúncios com poucas avaliações em direção à média global.

\subsection{Lacuna}

A literatura revisada documenta a compressão das notas, descreve os temas das avaliações do Airbnb e aplica modelos de linguagem à análise por aspecto. Faltam, em conjunto, quatro elementos que este trabalho reúne: a validação explícita da medida por aspecto com amostra ponderada pela população e com mais de um anotador, a comparação entre dois modelos de sentimento com regra de escolha fixada de antemão, a decomposição da nota global por aspecto e por tópico em um corpus multilíngue de um mercado brasileiro. A verificação de que o índice sintético construído a partir do texto não acrescenta poder explicativo à nota oficial.
